\section{Kombinatorika}
Notasikan $n!=n \times (n-1) \times (n-2) \times \dots \times 3 \times 2 \times 1$ (dibaca $n$ faktorial) dengan $1!=0!=1$.
\subsection{Aturan Pencacahan}
\input{High/SoalMath/Combinatorics/aturanPencacahan.tex}
\subsection{Permutasi}
Permutasi $k$ unsur dari $n$ unsur adalah (urutan diperhatikan)
$$_nP_k = P_k^n = \dfrac{n!}{(n-k)!}.$$
\input{High/SoalMath/Combinatorics/permutasi.tex}
\subsection{Kombinasi}
Kombinasi $k$ unsur dari $n$ unsur adalah (urutan tak diperhatikan)
$${n \choose k}=_nC_k = C_k^n = \dfrac{n!}{k!(n-k)!}.$$
\input{High/SoalMath/Combinatorics/kombinasi.tex}
\subsection{Permutasi Siklis}
$n$ objek ditaruh mengelilingi lingkaran maka banyak cara menyusunnya adalah
$$P_{siklis} =\dfrac{n!}{n} = (n-1)!$$
\input{High/SoalMath/Combinatorics/permutasiSiklis.tex}
\subsection{Stars and Bars}
Banyaknya solusi bulat non-negatif $(x_1,x_2,\dots,x_k)$ dari sistem persamaan $x_1+x_2+\dots+x_k=n$ adalah
$${n+k-1 \choose k-1}.$$
Banyaknya solusi bulat positif $(x_1,x_2,\dots,x_k)$ dari sistem persamaan $x_1+x_2+\dots+x_k=n$ adalah
$${n-1 \choose k-1}.$$
\input{High/SoalMath/Combinatorics/starsAndBars.tex}
\subsection{Prinsip Inklusi Eksklusi}
Prinsip Inklusi (memasukkan, dari kata inklusif) Eksklusi (mengeluarkan, khusus, dari kata eksklusif) atau yang sering disingkat PIE, pada dasarnya adalah konsep dari mengurangi "kelebihan hitung" atau menambahkan "kekurangan hitung". Contohnya adalah soal himpunan yang dinyatakan dalam rumus berikut
$$|A \cup B|=|A|+|B|-|A \cap B|.$$

Untuk tiga himpunan $A,B,C$ adalah
$$|A \cup B \cup C|=|A|+|B|+|C|-|A \cap B|-|A \cap C|-|B \cap C|+|A \cap B \cap C|.$$

dan seterusnya. Lebih lengkapnya boleh mengacu ke \href{https://brilliant.org/wiki/principle-of-inclusion-and-exclusion-pie/}{https://brilliant.org/wiki/principle-of-inclusion-and-exclusion-pie/}

\subsubsection{Derangement}
Teorema ini juga bisa disebut "teorema kado silang". Bunyi teorema ini:

Misalkan $n$ adalah bilangan bulat non-negatif. Kita sebut $!n$ atau $D_n$ sebagai derangement dari $n$ yaitu banyaknya permutasi $n$ elemen berbeda sedemikian sehingga tidak ada elemen yang menempati tempatnya semula.

\textbf{Versi yang tidak terlalu abstrak:} $!n$ adalah derangement dari $n$, dimana misalkan pada sebuah pesta ulang tahun, $n$ orang saling bertukar kado (awalnya semua orang mempunyai tepat satu kado) dimana setelah bertukar kado tidak ada orang yang mendapat kado dari dirinya sendiri. Banyak kemungkinan pertukaran kado ini adalah $!n$.

Rumus umum untuk menghitung derangement adalah
$$!n = n! \left(\dfrac{1}{0!}-\dfrac{1}{1!}+\dfrac{1}{2!}-\dfrac{1}{3!}+\dfrac{1}{4!}-\dfrac{1}{5!}+\dots+(-1)^n\dfrac{1}{n!}\right).$$
\input{High/SoalMath/Combinatorics/prinsipInklusiEksklusi.tex}
\subsection{Identitas Kombinatorika}
\begin{enumerate}
    \item  ${n \choose k} = {n \choose n-k}$ dengan $k,n \in \NN_0$ dan $k \le n$.
    \item (Identitas Pascal) Untuk $n,k \in \NN_0$ berlaku ${n \choose k} + {n \choose k+1} = {n+1 \choose k+1}$
    \item (Hockey Stick Identity) Untuk $n,r\in\mathbb{N}, n>r,\sum^n_{i=r}{i\choose r}={n+1\choose r+1}$.
    \item Untuk $n,k \in \NN_0$ berlaku ${n \choose 0}+{n \choose 1} + \dots + {n \choose n} = 2^{n}$.
    \item Untuk $n,k \in \NN_0$ berlaku ${n \choose 0}+{n \choose 2} + \dots + {n \choose 2\floor{\frac{n}{2}}} = 2^{n-1}$.
    \item (Vandermonde's Identity)  $\sum_{k=0}^r\binom mk\binom n{r-k}=\binom{m+n}r$
\end{enumerate}
\input{High/SoalMath/Combinatorics/identitasKombinatorika.tex}
\subsection{Binomial Newton}
Binomial Newton atau ekspansi/penjabaran binomial berfokus pada nilai koefisien setiap suku hasil penjabaran $(a+b)^n$.
$$(a+b)^n = {n \choose 0} a^nb^0 + {n \choose 1} a^{n-1}b^1+ \dots +{n \choose n}a^0b^n$$
\input{High/SoalMath/Combinatorics/binomialNewton.tex}
\subsection{Pigeon Hole Principle (PHP)}
Teorema yang dalam Bahasa Indonesia ini disebut dengan Teorema Sangkar Burung Merpati secara matematis berbunyi:
Jika ada $kn+1$ merpati dan $n$ sangkar, maka setidaknya ada satu sangkar yang berisi $k+1$ burung merpati.

Versi lebih simpelnya adalah: jika ada $n+1$ objek yang akan dibagi ke dalam $n$ buah kotak, maka setidaknya ada 1 kotak yang berisi 2 objek.

Contoh: \begin{itemize}
    \item Di dalam ruangan berisi 3 orang, pasti terdapat setidaknya 2 orang berjenis kelamin sama.
    \item Jika ada 367 orang di suatu sekolah, maka setidaknya ada dua orang diantara mereka yang tanggal lahirnya persis sama.
\end{itemize}
\input{High/SoalMath/Combinatorics/pigeonHolePrinciple.tex}
\subsection{Peluang}
Misalkan kita melempar sekeping koin, maka kegiatan ini disebut dengan percobaan. Hasil percobaan yang didapat biasanya adalah munculnya sisi gambar, $G$, atau munculnya sisi tulisan, $T$. Ruang contoh atau ruang sampel adalah himpunan dari \textbf{semua hasil percobaan yang mungkin} biasanya dilambangkan dengan $S$, yang dalam teori himpunan disebut dengan himpunan semesta. Pada percobaan melempar koin, ruang sampelnya adalah $\{G, T\}$ sedangkan pada percobaan melempar satu buah dadu, ruang sampelnya adalah $\{1, 2, 3, 4, 5, 6\}$. Jika $\{G, T\}$ adalah ruang sampel, maka anggota-anggota dari ruang sampel tersebut disebut titik contoh. Titik contoh dari $\{G, T\}$ adalah $G$ dan $T$. Pada percobaan melempar satu buah dadu seimbang, titik sampel yang didapat ada 6 yaitu 1, 2, 3, 4, 5, 6 sedangkan jika melempar dua buah dadu akan didapat 36 buah titik contoh, yaitu $(1, 1), (1, 2), (1, 3), \dots , (6, 6)$. 

\subsubsection{Formula Penghitungan Peluang}
Secara mudahnya, enghitung peluang bisa dengan pendekatan frekuensi, yaitu suatu percobaan yang dilakukan sebanyak $n$ kali, ternyata kejadian $A$ munculnya sebanyak $k$ kali, maka frekuensi nisbi/relatif kejadian $A$ sama dengan 
$$p(A)=\dfrac{k}{n}$$
Kalau $n$ semakin besar dan menuju tak terhingga maka nilai $p(A)$ akan cenderung konstan mendekati suatu nilai tertentu yang disebut dengan peluang munculnya kejadian $A$.
\input{High/SoalMath/Combinatorics/peluang.tex}
\subsection{Relasi Rekurensi}
Sering disebut dengan rekursif. Intinya adalah sebuah persamaan yang melibatkan barisan $a_1, a_2, \dots , a_n$ dimana untuk mendapatkan nilai $a_k$ membutuhkan suku-suku sebelumnya $a_{k-1}, a_{k-2}, \dots,$ atau $a_1$. 

Contoh paling terkenal dari persamaan rekursif adalah bilangan Fibonacci $0,1,1,2,3,5,8,13,21,\dots$ yang secara matematis didefinisikan sebagai berikut.
\begin{align*}
    F_0 &= 0, F_1 = 1\\
    F_n &= F_{n-1}+F_{n-2} \text{ untuk } n \ge 2
\end{align*}

atau yang lebih terkenal di ranah \textit{Computer Science} adalah permasalahan \textit{Tower of Hanoi} dengan persamaan rekursifnya didefinisikan sebagai berikut.
\begin{align*}
    T_1 &= 1 \\
    T_n &= 2T_{n-1}+1
\end{align*}

Untuk menyelesaikan soal relasi rekurensi, butuh manipulasi aljabar yang mumpuni sehingga tidak ada pendekatan eksplisit selain menggunakan persamaan karakteristik atau fungsi pembangkit (tidak dibahas disini) yang dijamin berhasil.

\subsubsection{Persamaan Karakteristik untuk Relasi Rekurensi Linear}
Persamaan karakteristik berikut berlaku untuk persamaan rekursif yang linear. Persamaan karakteristik berikut berguna untuk mengubah relasi rekurensi menjadi iteratif, atau persamaan berbentuk implisit. (Jadi, untuk persamaan yang bukan linear, sebagai contoh $a_n = a_{n-1}^2 + a_{n-2}$ tidak bisa dijamin selesai dengan persamaan karakteristik yang disajikan berikut).
Untuk persamaan rekursif
$$a_n = c_1a_{n-1}+c_2a_{n-2}+\dots+c_da_{n-d}$$
mempunyai persamaan karakteristik
$$x^d-c_1x^{d-1}-c_2x^{d-2}-\dots-c_dx^0=0$$

Sebagai contoh, rumus rekursif dari barisan Fibonacci di atas dapat diselesaikan menjadi 
$$x^{n}-x^{n-1}-x^{n-2}=0 \implies x^2-x-1=0$$
yang mempunyai dua akar, yaitu $x_1 = \dfrac{1+\sqrt{5}}{2}=\phi$ dan $x_2 = \dfrac{1-\sqrt{5}}{2}=1-\phi$ dimana $\phi$ adalah \textit{Golden Ratio}. Sadari bahwa setiap suku di barisan Fibonacci tersebut berbentuk $F_n = c_1x_1^n + c_2x_2^n$ (buktikan). Dengan substitusi $x_1$ dan $x_2$ serta pemilihan suku dari barisan Fibonacci (misal suku pertama dan kedua) maka akan ditemukan nilai $c_1$ dan $c_2$ sehingga pada akhirnya kita punya
$$F_n=\dfrac{\phi^n-(1-\phi)^n}{\sqrt{5}}$$
\input{High/SoalMath/Combinatorics/relasiRekurensi.tex}
